% Part: first-order-logic
% Chapter: sequent-calculus
% Section: quantifier-rules

\documentclass[../../../include/open-logic-section]{subfiles}

\begin{document}

\olfileid{fol}{seq}{qrl}

\olsection{পরিমাণসূচকের বিধি}

\subsection{$\lforall$-এর বিধি}

\begin{defish}
\Axiom$ !A(t), \Gamma \fCenter \Delta$
\RightLabel{\LeftR{\lforall}}
\UnaryInf$ \lforall[x][!A(x)],\Gamma \fCenter \Delta$
\DisplayProof
\hfill
\Axiom$ \Gamma \fCenter \Delta, !A(a) $
\RightLabel{\RightR{\lforall}}
\UnaryInf$ \Gamma \fCenter \Delta, \lforall[x][!A(x)]$
\DisplayProof
\end{defish}

\LeftR{\lforall}-এ $t$ একটি বদ্ধ পদ (অর্থাৎ চলরাশিবিহীন পদ)।
\RightR{\lforall}-এ $a$ এমন একটি !!a{constant}, যা
\RightR{\lforall} বিধির নিচের সিকোয়েন্টের কোথাও থাকতে পারবে না।
\RightR{\forall} অনুমানে $a$-কে \emph{আইগেনচল} বলা হয়।
\footnote{উপরের বিধিতে $a$ একটি !!a{constant} হলেও আমরা “আইগেনচল”
শব্দটি ব্যবহার করি। এর কারণ ঐতিহাসিক।}

\subsection{$\lexists$-এর বিধি}

\begin{defish}
\Axiom$ !A(a), \Gamma \fCenter \Delta $
\RightLabel{\LeftR{\lexists}}
\UnaryInf$ \lexists[x][!A(x)], \Gamma \fCenter \Delta$
\DisplayProof
\hfill
\Axiom$ \Gamma \fCenter \Delta, !A(t) $
\RightLabel{\RightR{\lexists}}
\UnaryInf$ \Gamma \fCenter \Delta, \lexists[x][!A(x)]$
\DisplayProof
\end{defish}

এখানেও $t$ একটি বদ্ধ পদ, আর $a$ এমন একটি !!a{constant}, যা
\LeftR{\lexists} বিধির নিচের সিকোয়েন্টে নেই। \LeftR{\lexists}
অনুমানে $a$-কে \emph{আইগেনচল} বলা হয়।

\RightR{\lforall} বা \LeftR{\lexists} অনুমানের নিচের সিকোয়েন্টে
আইগেনচল না থাকার শর্তটিকে \emph{আইগেনচল-শর্ত} বলা হয়।

\begin{explain}
মনে করুন, $!A$ কোনো !!a{formula} এবং তাতে !!{variable}~$x$ মুক্ত হলে
আমরা~$!A(x)$ লিখে তা বোঝাই। একই প্রসঙ্গে $!A(t)$ হলো
সংক্ষেপে~$\Subst{!A}{t}{x}$। তাই \RightR{\lexists} বিধিটি এভাবেও
লেখা যায়:
\begin{prooftree}
  \Axiom$\Gamma \fCenter \Delta, \Subst{!A}{t}{x}$
  \RightLabel{\RightR{\lexists}}
  \UnaryInf$\Gamma \fCenter \Delta, \lexists[x][!A]$
\end{prooftree}
লক্ষ করুন,~$!A$-তে $t$ আগে থেকেই থাকতে পারে; যেমন, $!A$ হতে পারে
~$\Atom{\Obj P}{t,x}$। অতএব~$\Gamma \Sequent \Delta,
\Atom{\Obj P}{t,t}$ থেকে $\Gamma \Sequent \Delta,
\lexists[x][\Atom{\Obj P}{t,x}]$ অনুমান করা
\RightR{\lexists}-এর সঠিক প্রয়োগ—পূর্বধারণায় $t$-এর এক বা একাধিক,
কিন্তু আবশ্যিকভাবে সব নয়, এমন উপস্থিতিকে বদ্ধ !!{variable}~$x$ দিয়ে
“প্রতিস্থাপন” করা যায়। কিন্তু \RightR{\lforall} ও~\LeftR{\lexists}-এর
আইগেনচল-শর্ত অনুযায়ী~$!A$-তে !!{constant}~$a$ থাকতে পারে না। তাই
$\Gamma \Sequent \Delta, \Atom{\Obj P}{a,a}$ থেকে
$\Gamma \Sequent \Delta, \lforall[x][\Atom{\Obj P}{a,x}]$
~$\RightR{\lforall}$ ব্যবহার করে সঠিকভাবে অনুমান করা যায় না।
\end{explain}

\begin{explain}
\RightR{\lexists} ও \LeftR{\lforall}-এ, আগে বলা বদ্ধতার শর্ত ছাড়া,
$t$-এর উপর অন্য কোনো বিধিনিষেধ নেই। অন্যদিকে, \LeftR{\lexists} ও
\RightR{\lforall} বিধির আইগেনচল-শর্ত অনুযায়ী উপরের সিকোয়েন্টে
$!A(a)$-এর বাইরে !!{constant}~$a$ কোথাও থাকতে পারে না। পদ্ধতিটি বিশুদ্ধ,
অর্থাৎ শুধু বৈধ সিকোয়েন্টই !!{derive}s করে—এটি নিশ্চিত করার জন্য শর্তটি
আবশ্যক। এই শর্ত না থাকলে নিচের অনুমানগুলি অনুমোদিত হতো:
\begin{prooftree}
  \Axiom$!A(a) \fCenter !A(a)$
  \RightLabel{*\LeftR{\lexists}}
  \UnaryInf$\lexists[x][!A(x)] \fCenter !A(a)$
  \RightLabel{\RightR{\lforall}}
  \UnaryInf$\lexists[x][!A(x)] \fCenter \lforall[x][!A(x)]$
  \DisplayProof\bottomAlignProof
  \qquad
  \Axiom$!A(a) \fCenter !A(a)$
  \RightLabel{*\RightR{\lforall}}
  \UnaryInf$!A(a) \fCenter \lforall[x][!A(x)]$
  \RightLabel{\LeftR{\lexists}}
  \UnaryInf$\lexists[x][!A(x)] \fCenter \lforall[x][!A(x)]$
\end{prooftree}
কিন্তু $\lexists[x][!A(x)] \Sequent \lforall[x][!A(x)]$ বৈধ নয়।
\end{explain}

\end{document}
